Before we state our formal results in \secref{finite_dim_clts}, we will present
some high-level intuition to motivate the IJ covariance estimate, as well as our
data-dependent Bayesian central limit theorem of \secref{finite_dim_clts}.
%   Additionally, we
% will now motivate why the $O_p$-style bounds of existing Bayesian central limit
% theorems are insufficient to prove the consistency of the IJ covariance
% estimate.

For simplicity let us for the moment take the quantity of interest,
$g(\theta)$, to simply be $\g(\theta) = \theta$.
% Our
% formal results will be stated for generic $\g(\theta)$ and will follow
% from this
% the delta method for general differentiable $g(\theta)$.
%
% Classical CLT results, and the target variance.
Following \citet{kleijn:2012:bvm}, for example, we might informally expect the
following central limit theorem results to hold:
%
\begin{align}
%
N \cov{\post}{\theta} =&{} \info^{-1} + O_p(N^{-1})
    &\textrm{Bayesian CLT (BCLT)} \nonumber \\
% \sqrt{N}\left(\expect{\post}{\theta} - \thetahat \right) \plim& 0
%     &\textrm{(Bayesian CLT)}\\
\expect{\post}{\theta} =&{} \thetahat + O_p(N^{-1}).
    &\textrm{Bayesian CLT (BCLT)} \nonumber \\
\sqrt{N}\left(\thetahat - \thetatrue\right) \dlim&{}
    \normdist\left(\cdot | 0, \info^{-1} \scorecov \info^{-1} \right).
    &\textrm{Frequentist CLT (FCLT)}
    \eqlabel{fclt_bclt}
%
\end{align}
%
Here and throughout the paper, we somewhat loosely use the term ``Bayesian CLT''
(BCLT) to refer to results concerning the asymptotic properties of the posterior
measure, and ``Frequentist CLT'' (FCLT) for results concerning properties under
the data sampling measure.  Both BCLT and FCLT results are typically proven
using frequentist tools.

Given the work of \citet{kleijn:2012:bvm}, we know that the limiting covariance we
are trying to estimate is the sandwich covariance matrix, $\gcovtrue =
\info^{-1} \scorecov \info^{-1}$. Since $N \cov{\post}{\theta}$ converges to
$\info^{-1}$, but $\info^{-1} \ne \info^{-1} \scorecov \info^{-1}$
in general unless the model is correctly specified, $\gcovtrue$ cannot,
in general, be accurately approximated by the Bayesian posterior covariance.
%
% We could just use the plug-in estimator if it were available.
A natural estimator of the sandwich covariance is the plug-in estimator $\gcovhat
:= \infohat^{-1} \scorecovhat \infohat^{-1}$, which consistently estimates
$\gcovtrue$.  However, computing $\gcovhat$ requires the MAP $\thetahat$, which is not
typically computable directly from MCMC samples.

% The IJ covariance
% estimator is motivated by circumstances when it is not easy or desirable to
% compute $\thetahat$, but when one can compute or approximate the posterior
% expectations required to compute the IJ covariance estimator $\gcovij$.

% Taylor expanding is the trick
The reason that $\gcovij$ might be expected to approximate $\gcovtrue$
asymptotically can be seen by applying the delta method and the
CLT results of \eqref{fclt_bclt} to the
influence score $\infl_n$ defined in \eqref{ij_terms}.  Specifically,
for any particular $n$, we expect that
%
\begin{align}\eqlabel{infl_heuristic}
    \infl_n = N \cov{\post}{\theta, \ell(\x_n \vert \theta)} =
    \infohat^{-1} \ellgrad{1}(\x_n \vert \thetahat) + O_p(N^{-1}).
\end{align}
%
% (See \appref{bclt_uses} for a more detailed derivation.)
If \eqref{infl_heuristic} holds sufficiently uniformly in $n$
(in a sense that we analyze precisely below), then \eqref{infl_heuristic}
implies the asymptotic equivalence of $\gcovij$ and $\gcovtrue$:
%
\begin{align}
%
% \bar{\infl} \approx{}& \info^{-1} \meann \ellgrad{1}(\x_n \vert \thetahat) = 0
% \quad\quad \textrm{and} \nonumber\\
\infl(\x_n) \approx{}& \infohat^{-1} \ellgrad{1}(\x_n \vert \thetahat)
\quad\Rightarrow\quad
\gcovij ={} \meann \infl_n \infl_n^T - \bar{\infl} \, \bar{\infl}^T
% \nonumber\\
% \approx{}&
%     \infohat^{-1} \left( \meann
%         \ellgrad{1}(\x_n \vert \thetahat)
%         \ellgrad{1}(\x_n \vert \thetahat)^T \right)
%     \infohat^{-1}
% \nonumber\\
\approx{} \infohat^{-1} \scorecovhat \infohat^{-1}, \eqlabel{ij_heuristic}
% = \gcovhat \plim \gcovtrue.
%
\end{align}
%
and we know that $\infohat^{-1} \scorecovhat \infohat^{-1} = \gcovhat \plim
\gcovtrue$.


However, in general, the $\ordp{N^{-1}}$ term in \eqref{infl_heuristic}
depends on $\x_n$, and the IJ covariance depends on {\em every} $\infl(\x_n): n \in
[N]$, some of which may be pathologically behaved, especially under
misspecification.
%
A rigorous analysis of the IJ covariance requires
understanding whether the approximation in \eqref{infl_heuristic} holds
sufficiently strongly for us to apply the reasoning of \eqref{ij_heuristic}.
In particular, we must establish that we can safely sum the squares of the $\ordp{N^{-1}}$
terms of \eqref{infl_heuristic}.  It is to this technical task that we
now turn.

% Furthermore, it is not reasonable to assume that
% the BCLT residual is uniformly small for all datapoints $\x_n$. After all, the
% most extreme datapoints in a growing dataset might be quite extreme indeed.


%%%%%%%%%%%%%%%%%%%%%%%%%%%%%%%%%%%%%%%%%%%%%%%%%%%%
